\documentclass[12pt]{article}
\usepackage[utf8]{inputenc}
\usepackage[T1]{fontenc}
\usepackage[margin=1in]{geometry}
\usepackage{setspace}
\usepackage{xcolor}
\usepackage[colorlinks=true, urlcolor=blue!60!black, linkcolor=blue!60!black]{hyperref}
\setlength{\parindent}{2em}
\setlength{\parskip}{1.25em}
\renewcommand{\baselinestretch}{1.0}
\begin{document}
\begin{center}
{\LARGE \textit{Brandenburg v. Ohio}, 395 U.S. 444 (1969)\par}
\vspace{0.8em}
{\large Case Brief --- Four-Page Edition\par}
\end{center}
\vspace{1em}

\section{Memory Jogger}

A Ku Klux Klan leader convicted under Ohio's criminal syndicalism law for a filmed rally speech hinting at ``revengeance'' won reversal: a unanimous per curiam Court held that the State may punish advocacy of force or law violation only when it is directed to inciting imminent lawless action and is likely to produce it, and overruled \textit{Whitney v.\ California} (\href{https://supreme.justia.com/cases/federal/us/395/444/}{Justia}).

\section{Detailed Case Facts}

Clarence Brandenburg, a Klan leader, invited a Cincinnati television reporter to cover a Klan ``rally'' on a Hamilton County farm. Only Klan members and the news crew attended. One film showed twelve hooded figures, some armed, around a burning cross, uttering phrases derogatory of Black people and Jews. Brandenburg, unarmed and in regalia, said: ``We're not a revengent organization, but if our President, our Congress, our Supreme Court, continues to suppress the white, Caucasian race, it's possible that there might have to be some revengeance taken.'' He announced a July Fourth march on Congress ``four hundred thousand strong.'' A second film showed a similar speech urging that Black people be returned to Africa and Jews to Israel. Ohio introduced the films and items shown in them, including firearms and a red hood.

\section{Procedural History}

Brandenburg was convicted under Ohio's Criminal Syndicalism Statute for advocating ``the duty, necessity, or propriety of crime, sabotage, violence, or unlawful methods of terrorism as a means of accomplishing industrial or political reform'' and for assembling with a group formed to advocate it. He was fined \$1,000 and sentenced to one to ten years. An intermediate court affirmed, and the Supreme Court of Ohio dismissed his appeal for lack of a ``substantial constitutional question.'' The U.S. Supreme Court reversed on June 9, 1969.

\section{Judicial Votes}

Reversed, \textbf{8--0}, \emph{per curiam}.
\begin{itemize}
\item \textbf{Per curiam:} Chief Justice Warren and Justices Black, Douglas, Harlan, Brennan, Stewart, White, and Marshall (one seat vacant after Justice Fortas resigned).
\item \textbf{Concurring:} Justice Black; Justice Douglas.
\end{itemize}

\section{Holding}

Ohio's statute is invalid on its face and as applied. The First and Fourteenth Amendments forbid a State to proscribe advocacy of force or law violation except where it is (1) directed to inciting or producing imminent lawless action and (2) likely to incite or produce such action. The same distinction governs restrictions on peaceable assembly. \textit{Whitney v.\ California} (1927) is overruled.

\section{Analysis of Opinions}

\textbf{Per curiam.} Ohio's 1919 law was one of many criminal syndicalism statutes of 1917--1920. \textit{Whitney} had upheld California's similar statute, but it had been ``thoroughly discredited by later decisions,'' including \textit{Dennis} and \textit{Yates}. The Court distilled a two-part test: an intent and immediacy requirement and an objective likelihood requirement. Quoting \textit{Noto v.\ United States}, it distinguished ``mere abstract teaching'' of the propriety of violence from ``preparing a group for violent action and steeling it to such action.'' Ohio's statute punished advocacy, teaching, and assembly without that distinction, and neither the indictment nor the jury instructions, which tracked the statute, narrowed it.

\textbf{Black (concurring).} He joined on the understanding that citing \textit{Dennis} did not endorse its ``clear and present danger'' doctrine, which should have no place in First Amendment law.

\textbf{Douglas (concurring).} Tracing ``clear and present danger'' from Holmes's \textit{Schenck} opinion to its Cold War distortion in \textit{Dennis}, he found it manipulable and incompatible with the First Amendment in peacetime, and drew a categorical line between ``ideas and overt acts.''

\textbf{Later cases.} \textit{Hess v.\ Indiana} (1973) held that advocacy of illegal action at ``some indefinite future time'' fails the imminence prong; \textit{NAACP v.\ Claiborne Hardware} (1982) applied \textit{Brandenburg} to civil liability; and \textit{Counterman v.\ Colorado} (2023) confirmed that incitement requires specific intent.

\section{Examples: Future Applications}

\begin{enumerate}
\item \textbf{Same-Side: Revolutionary Pamphlet.} A group argues that armed revolution will someday be morally necessary, naming no target or date. This is protected ``mere abstract teaching,'' and a statute without an imminence requirement is facially invalid.
\item \textbf{Same-Side: Boycott Warning.} An organizer warns that boycott violators will regret it; weeks later unknown persons vandalize a car. Under \textit{Claiborne Hardware}, no civil liability.
\item \textbf{Opposite-Side: Bullhorn Command.} Outside a besieged building, a leader tells an armed crowd of 200, ``The doors are unlocked right now---rush the building and drag the councilmembers out!'' Lawless action is directed, imminent, and likely.
\item \textbf{Opposite-Side: Solicitation of Arson.} A private offer of \$5,000 to burn a warehouse that night is solicitation, unprotected under \textit{United States v.\ Williams}.
\item \textbf{Fence-Sitter: Scheduled Livestream Raid.} An influencer tells 500,000 followers to bring crowbars to a courthouse at noon tomorrow. Under \textit{Hess}, imminence is a substantial hurdle, though specificity bears on intent and likelihood.
\end{enumerate}

\section{Critique}

\textbf{Doctrinal revisionism and vagueness.} The per curiam claimed to synthesize \textit{Dennis} and \textit{Yates}, though \textit{Dennis} used a gravity-discounted-by-improbability balancing test without strict imminence. And \textit{Brandenburg} never defined how immediate ``imminent'' must be, or whether ``lawless action'' includes minor misdemeanors such as trespass.

\textbf{Digital threats.} Critics argue that the paradigm of a speaker addressing a live crowd fits poorly with asynchronous online radicalization, viral stochastic incitement, and instructional manuals for crimes.

\section{Key Quotations}

\begin{quote}
``[T]he constitutional guarantees of free speech and free press do not permit a State to forbid or proscribe advocacy of the use of force or of law violation except where such advocacy is directed to inciting or producing imminent lawless action and is likely to incite or produce such action.''
\end{quote}
\noindent --- \emph{Per curiam}, 395 U.S.\ at 447.

\begin{quote}
``Accordingly, we are here confronted with a statute which, by its own words and as applied, purports to punish mere advocacy and to forbid, on pain of criminal punishment, assembly with others merely to advocate the described type of action. Such a statute falls within the condemnation of the First and Fourteenth Amendments. The contrary teaching of \emph{Whitney v.\ California}, \emph{supra}, cannot be supported, and that decision is therefore overruled.''
\end{quote}
\noindent --- \emph{Per curiam}, 395 U.S.\ at 449.

\begin{quote}
``The line between what is permissible and not subject to control and what may be made impermissible and subject to regulation is the line between ideas and overt acts. The example usually given by those who would punish speech is the case of one who falsely shouts fire in a crowded theatre. This is, however, a classic case where speech is brigaded with action.''
\end{quote}
\noindent --- Justice Douglas, concurring, 395 U.S.\ at 456--57.

\end{document}
